\begin{figure}[ht]
  \centering
  \begin{tikzpicture}
    \begin{axis}[paperaxis,
      ybar, bar width=14pt,
      ylabel={Réponse à l'impôt (valeur absolue)},
      symbolic x coords={distribution,retention,fondspropres},
      xtick=data,
      xticklabels={Distribution, Réserves, Fonds propres},
      ymin=0,
      enlarge x limits=0.28,
      legend pos=north east, legend cell align=left,
      nodes near coords, nodes near coords style={font=\scriptsize},
      every node near coord/.append style={/pgf/number format/.cd, use comma, fixed, fixed zerofill, precision=2},
    ]
      % La hausse de 2013 en orange, la baisse de 2018 en vert : meme code que partout.
      \addplot[fill=hausse2013, draw=hausse2013!60!black]
        table[col sep=comma, x=marge, y=choc2013] {figures/cliquet.csv};
      \addlegendentry{Choc de hausse, 2013}
      \addplot[fill=baisse2018, draw=baisse2018!60!black]
        table[col sep=comma, x=marge, y=choc2018] {figures/cliquet.csv};
      \addlegendentry{Choc de baisse, 2018}
    \end{axis}
  \end{tikzpicture}
  \caption{Le cliquet fiscal : la réponse à la hausse de 2013 dépasse la réponse à la baisse de 2018}
  \label{fig:cliquet}
  \fignote{Ampleur de la réponse à chaque choc, en valeur absolue, sur trois marges. La distribution
  est la semi-élasticité des dividendes des ménages, sur base lissée des trois années antérieures au
  choc ; la rétention et les fonds propres sont les coefficients de double différence sur les parts
  de l'actif, au niveau des sociétés. Par convention de couleur, l'orange est la hausse d'impôt de
  2013 et le vert la baisse de 2018. Sur les trois marges, la réponse à la hausse de 2013 dépasse la
  réponse à la baisse de 2018 : le capital part plus facilement qu'il ne revient. La semi-élasticité
  de la distribution est négative en 2013 et positive en 2018 ; seules les amplitudes sont tracées.}
\end{figure}
